\section*{Bab \ref{ch:g3:separating-mixtures} --- Memisahkan Campuran}

\begin{solution}{exo:g3:separating-mixtures:1}
\omterm{def:g3:separating-mixtures:sieve}{Pengayakan}: butiran beras lebih kecil daripada kacang polong dan jatuh
melewati ayakan yang lubangnya pas. (Memilah dengan tangan juga
bisa, tetapi lambat.)
\end{solution}

\begin{solution}{exo:g3:separating-mixtures:2}
Tanahnya perlahan-lahan turun dan membentuk lapisan di dasar; air di
atasnya menjadi lebih jernih. Inilah \omterm{def:g3:separating-mixtures:decant}{pengendapan}.
\end{solution}

\begin{solution}{exo:g3:separating-mixtures:3}
\omterm{def:g3:separating-mixtures:filter}{Filtrat} adalah cairan yang lolos melewati saringan. Ketika orang
membuat kopi, \omterm{def:g3:separating-mixtures:filter}{filtratnya} adalah kopi di dalam teko; ampasnya tertahan
di kertas.
\end{solution}

\begin{solution}{exo:g3:separating-mixtures:4}
Ya. Garamnya \omterm{def:g2:mixing-and-dissolving:dissolve}{larut}, dan saringan tidak menahan apa yang \omterm{def:g2:mixing-and-dissolving:dissolve}{larut}: garam
itu lolos bersama airnya.
\end{solution}

\begin{solution}{exo:g3:separating-mixtures:5}
Dengan \omterm{def:g3:separating-mixtures:evaporate}{menguapkan} airnya: di bawah matahari, atau dipanaskan perlahan,
airnya pergi ke udara dan gulanya tertinggal di cawan.
\end{solution}

\begin{solution}{exo:g3:separating-mixtures:6}
Kerikil tertahan di ayakan. Butiran pasir lebih kecil daripada
lubangnya, jadi butiran itu jatuh melewatinya.
\end{solution}

\begin{solution}{exo:g3:separating-mixtures:7}
(a) \omterm{def:g3:separating-mixtures:sieve}{Pengayakan}. (b) \omterm{def:g3:separating-mixtures:filter}{Penyaringan} (atau \omterm{def:g3:separating-mixtures:decant}{pengendapan} lalu \omterm{def:g3:separating-mixtures:decant}{dekantasi}).
(c) \omterm{def:g3:separating-mixtures:evaporate}{Penguapan}.
\end{solution}

\begin{solution}{exo:g3:separating-mixtures:8}
Empat langkah: jeruji, bak pengendap, saringan pasir, pembasmian kuman.
Bak pengendap membuang lumpurnya.
\end{solution}

\begin{solution}{exo:g3:separating-mixtures:9}
Saring; ambil pasir dari kertas saring; \omterm{def:g3:separating-mixtures:evaporate}{uapkan} \omterm{def:g3:separating-mixtures:filter}{filtratnya} untuk
mendapatkan garam.
\end{solution}

\begin{solution}{exo:g3:separating-mixtures:10}
Tidak. Garam dalam air laut \omterm{def:g2:mixing-and-dissolving:dissolve}{larut}, dan saringan tidak menahan apa yang
\omterm{def:g2:mixing-and-dissolving:dissolve}{larut}: \omterm{def:g3:separating-mixtures:filter}{filtratnya} tetap air asin.
\end{solution}

\begin{solution}{exo:g3:separating-mixtures:11}
Menuang tehnya perlahan setelah daunnya mengendap (\omterm{def:g3:separating-mixtures:decant}{dekantasi}), atau
menuangnya melalui saringan teh atau kertas saring (\omterm{def:g3:separating-mixtures:sieve}{pengayakan} atau
\omterm{def:g3:separating-mixtures:filter}{penyaringan}).
\end{solution}
